% ch00 导读说明 (原创)
\chapter*{导读说明}
\addcontentsline{toc}{chapter}{导读说明}

本册是 Belal E. Baaquie所著 \emph{Quantum Finance: Path Integrals and
Hamiltonians for Options and Interest Rates}（Cambridge University Press,
2004）的\textbf{中文精读导读}，供持有原书的读者对照阅读。它\textbf{不是译本}：
全书以讲解体重新组织——每节先说明"这一节要解决什么问题"，再按逻辑顺序
呈现模型设定与关键公式（编号与原书一致，形如 (N.M)），最后用自己的话
解释推导主线与金融/物理含义。原书的完整行文、例题展开与文献细节仍以
原书为准；本册的作用是把全书的数学骨架提出来，让读者知道每一步"为什么
这么走"，并在回头精读原书时有所依凭。

\section*{全书结构一览}

原书论证的主线是一条不断"升级物理工具"的路线：

\begin{enumerate}[itemsep=0.2em]
\item 第 1 章给出全书 synopsis：把期权定价表述为倒向随机演化的（虚时）
      量子力学问题，把利率曲线表述为一维"量子场"。
\item 第 2--3 章是金融学背景：有效市场、随机游走、无套利与鞅、远期/
      期货/期权、Black--Scholes 方程与 Merton--Garman 方程（随机波动率）。
\item 第 4--5 章把股票期权改写成量子力学语言：定价核 = 传播子，
      Black--Scholes 哈密顿量、路径积分表示、障碍期权 = 势垒问题。
\item 第 6 章把同样的机器用于即期利率模型（Vasicek 等），
      引入正演/倒向哈密顿量与 Fokker--Planck 哈密顿量。
\item 第 7--10 章是本书核心贡献：把整条远期利率曲线 $f(t,x)$ 看作
      二维量子场（$t$ 时间、$x$ 未来时间），构造场论作用量、传播子
      （刚性/刚度）、鞅条件与风险中性测度的场论形式，并用欧洲美元
      期货数据做实证检验（第 8 章），推广到国债衍生品定价与对冲
      （第 9 章）与哈密顿量形式（第 10 章）。
\item 附录 A 补数学工具：泛函高斯积分、白噪声、朗之万方程与
      金融学基本定理。
\end{enumerate}

\section*{使用约定}

\begin{itemize}[itemsep=0.2em]
\item \textbf{公式编号}与原书一致："式 (7.38)"即原书第 7 章第 38 式；
      原书无编号的推导行在本册中不加编号。个别推导步骤从略处均注明。
\item \textbf{图表}不复制：以"原书图 N.M / 原书表 N.M"指路，请对照原书。
      提取的原书插图片文件随本工程存于 \texttt{figures/} 文件夹
      （\texttt{fig\_N\_M.png}），仅供个人对照查阅。
\item \textbf{术语}首次出现括注英文；全书术语与记号见卷末
      "术语中英对照表"与"记号总表"。
\item \textbf{记号}完全沿用原书：股票价格 $S(t)$、对数价格
      $x(t)=\ln S(t)$、即期利率 $r(t)$、远期利率 $f(t,x)$、
      零息国债 $P(t,T)$、距到期时间 $\tau=T-t$、远期利率速度场
      $\mathcal{A}(t,x)$ 等；花体字母
      $\mathcal{A},\mathcal{B},\mathcal{F},\mathcal{H},\mathcal{L},
      \mathcal{V},\mathcal{O},\mathcal{P}$ 保留原书含义。
\item 虚时（Euclidean）时间与金融时间的对应是全书反复出现的主题：
      期权定价用倒向（backward）演化，国债定价用正演（forward）
      演化，两者由 Feynman--Kac 公式联系。
\end{itemize}

\section*{阅读建议}

熟悉量子力学与路径积分的读者：可直接从第 4 章进入，把第 2.5 节
（无套利与风险中性测度）当作金融侧的唯一"公理"补读；第 7 章的
场论构造是全书技术核心。熟悉金融的读者：第 2--3 章可快速通过，
注意第 4 章起"哈密顿量 $H = -\frac{\sigma^2}{2}\partial_x^2 +
\frac{\sigma^2}{2}\partial_x + r$ 与薛定谔形式"的对应，以及第 5 章
作用量 $S = \int L\,dt$ 与定价核 $p \sim e^{-S}$ 的路径积分表示。
第 8 章的实证部分（相关结构、刚性/刚度参数拟合、心理未来时间）
是理解第 7 章模型为何取该形式的关键证据。
